\hypertarget{mathematical-formulation-of-the-fish-metacommunity-model}{%
\section{Mathematical Formulation of the Fish Metacommunity
Model}\label{mathematical-formulation-of-the-fish-metacommunity-model}}

This document provides a detailed description of the spatially-explicit
modeling framework used to simulate fish metacommunity dynamics in the
Guadalquivir River basin. The model is implemented as a system of
Ordinary Differential Equations (ODEs) that integrate local population
dynamics, environmental filtering, and 3D dendritic dispersal.

\hypertarget{general-model-structure}{%
\subsection{1. General Model Structure}\label{general-model-structure}}

The rate of change for the population size \(N_{i,s}\) of species \(s\)
at site \(i\) is defined by the sum of local dynamics and net spatial
flux:

\[\frac{dN_{i,s}}{dt} = \underbrace{f_{local}(N_{i,s}, \mathbf{N}_i, E_i)}_{\text{Local Dynamics}} + \underbrace{\sum_{j \in \text{neighbors}} (m_{ji} N_{j,s} - m_{ij} N_{i,s})}_{\text{Spatial Flux}}\]

where: - \(N_{i,s}\) is the population density of species \(s\) at site
\(i\). - \(\mathbf{N}_i\) is the vector of all species populations at
site \(i\). - \(E_i\) represents the local environmental conditions. -
\(m_{ij}\) is the dispersal rate from site \(i\) to site \(j\).

\hypertarget{local-dynamics-and-environmental-filtering}{%
\subsection{2. Local Dynamics and Environmental
Filtering}\label{local-dynamics-and-environmental-filtering}}

The local component follows a Generalized Lotka-Volterra (GLV)
structure, modified by environmental suitability filters.

\hypertarget{population-growth-and-biotic-interactions}{%
\subsubsection{2.1. Population Growth and Biotic
Interactions}\label{population-growth-and-biotic-interactions}}

The local dynamics are expressed as:

\[f_{local} = N_{i,s} \left( r_{i,s}^{eff} - \sum_{j=1}^{S} \alpha_{sj} N_{i,j} \right)\]

where: - \(r_{i,s}^{eff}\) is the effective intrinsic growth rate. -
\(\alpha_{sj}\) is the interaction coefficient representing the effect
of species \(j\) on species \(s\) (asymmetric interaction matrix).

\hypertarget{environmental-filtering}{%
\subsubsection{2.2. Environmental
Filtering}\label{environmental-filtering}}

The effective growth rate \(r_{i,s}^{eff}\) is determined by the base
growth rate \(r_{s}\) and a site-specific suitability factor
\(W_{i,s}\):

\[r_{i,s}^{eff} = r_{s} \cdot W_{i,s}\]

The suitability factor \(W_{i,s}\) incorporates thermal tolerance and
habitat quality:

\[W_{i,s} = \text{ThermalFilter}(T_i, \text{opt}_s, \sigma_s) \cdot h_i\]

\hypertarget{thermal-filter}{%
\paragraph{Thermal Filter}\label{thermal-filter}}

We use a Gaussian-shaped ``thermal window'' to model species-specific
temperature tolerances:

\[\text{ThermalFilter}(T_i, \text{opt}_s, \sigma_s) = \exp\left( -\frac{(T_i - \text{opt}_s)^2}{2\sigma_s^2} \right)\]

where: - \(T_i\) is the temperature at site \(i\). - \(\text{opt}_s\) is
the optimal temperature for species \(s\). - \(\sigma_s\) is the thermal
tolerance (standard deviation) of species \(s\).

\hypertarget{habitat-suitability}{%
\paragraph{Habitat Suitability}\label{habitat-suitability}}

\begin{itemize}
\tightlist
\item
  \(h_i\) is the habitat suitability index at site \(i\), which accounts
  for anthropogenic stressors (e.g., conductivity), river width, and
  mesohabitat types.
\end{itemize}

\hypertarget{spatial-flux-3d-dendritic-dispersal}{%
\subsection{3. Spatial Flux (3D Dendritic
Dispersal)}\label{spatial-flux-3d-dendritic-dispersal}}

The spatial component models the movement of individuals through the
river network, accounting for hydrological distance, elevation changes,
and physical barriers.

\hypertarget{dispersal-rate-formulation}{%
\subsubsection{3.1. Dispersal Rate
Formulation}\label{dispersal-rate-formulation}}

The dispersal rate \(m_{ij}\) from site \(i\) to site \(j\) is defined
as:

\[m_{ij} = m \cdot \left( \frac{1}{d_{ij}} \right) \cdot x_{ij} \cdot p_{ij}\]

where: - \(m\) is the base dispersal intensity. - \(d_{ij}\) is the
hydrological distance between site \(i\) and \(j\). - \(x_{ij}\) is the
3D connectivity factor (elevation cost). - \(p_{ij}\) is the dam
passability factor (ranging from 0.0 to 1.0).

\hypertarget{elevation-cost-upstream-vs.-downstream}{%
\subsubsection{3.2. Elevation Cost (Upstream
vs.~Downstream)}\label{elevation-cost-upstream-vs.-downstream}}

The cost of movement depends on the relative elevation of the sites:

\begin{itemize}
\tightlist
\item
  \textbf{Downstream or Equal Elevation (\(e_j \le e_i\)):}
  \[x_{ij} = 1.0\]
\item
  \textbf{Upstream (\(e_j > e_i\)):}
  \[x_{ij} = \frac{1}{1 + c(e_j - e_i)}\]
\end{itemize}

where: - \(e_i\) and \(e_j\) are the elevations of the source and
destination sites, respectively. - \(c\) is the upstream migration cost
constant.

\hypertarget{implementation-details}{%
\subsection{4. Implementation Details}\label{implementation-details}}

\begin{itemize}
\tightlist
\item
  \textbf{State Vector:} The system state \(u\) is a matrix of size
  \(N_{sites} \times N_{species}\).
\item
  \textbf{Sparse Matrix Optimization:} To handle large-scale river
  networks efficiently, the dispersal rates are pre-calculated into a
  sparse matrix \(M\), where \(M_{ji}\) represents the rate from \(i\)
  to \(j\).
\item
  \textbf{Numerical Integration:} Given the potential stiffness arising
  from the coupling of fast dispersal and slower local dynamics, the
  system is solved using stiff-aware solvers (e.g., \texttt{Rodas5} or
  \texttt{TRBDF2}).
\end{itemize}

\hypertarget{model-assumptions}{%
\subsection{5. Model Assumptions}\label{model-assumptions}}

\begin{enumerate}
\def\labelenumi{\arabic{enumi}.}
\tightlist
\item
  \textbf{Constant Interaction Matrix:} Biotic interaction coefficients
  (\(\alpha_{sj}\)) are assumed to be constant across the entire basin,
  though the realized interaction depends on local densities.
\item
  \textbf{Passive/Active Dispersal Balance:} Dispersal is modeled as a
  density-dependent flux proportional to the population at the source
  site.
\item
  \textbf{Gaussian Thermal Niche:} Species performance is assumed to
  decay symmetrically as temperatures deviate from the optimum.
\item
  \textbf{Dam Passability:} Dams act as semi-permeable filters that
  reduce the effective dispersal rate between connected nodes.
\end{enumerate}
